\documentclass[12pt]{article}
\usepackage{fullpage}
\usepackage[T1]{fontenc}
\usepackage[utf8]{inputenc}
\usepackage{lmodern}
\usepackage{microtype}
\usepackage{amsmath,amssymb,amsthm}
\usepackage{hyperref}

\newtheorem{theorem}{Theorem}
\newtheorem{lemma}{Lemma}
\newtheorem{proposition}{Proposition}
\theoremstyle{remark}
\newtheorem*{remark}{Remark}

\DeclareMathOperator{\ran}{ran}
\newcommand{\spec}{\sigma}

\begin{document}

\title{Singular values of finite Hilbert matrices converge to $\pi$}
\author{}
\date{}
\maketitle

\begin{theorem}\label{thm:main}
Let $H_n=\bigl(\tfrac{1}{i+j-1}\bigr)_{i,j=1}^{n}$ be the $n\times n$ Hilbert matrix, and for
fixed $i\ge 1$ and $n\ge i$ let $\sigma_i(H_n)$ denote its $i$-th largest eigenvalue
(equivalently, since $H_n$ is symmetric positive definite, its $i$-th largest singular value).
Then for \emph{every} fixed index $i\ge 1$,
\[
\lim_{n\to\infty}\sigma_i(H_n)=\pi .
\]
In particular the proposed equality \emph{holds} (it does not fail).
Moreover, for each fixed $i$ the sequence $\bigl(\sigma_i(H_n)\bigr)_{n\ge i}$ is nondecreasing
and bounded above by $\pi$, with $\sigma_i(H_n)<\pi$ for every finite $n$, and the convergence is
only logarithmically fast: $0<\pi-\sigma_i(H_n)=O\!\left(1/\log n\right)$.
\end{theorem}

The proof rests on the spectral theory of the \emph{infinite} Hilbert matrix, viewed as a
bounded self-adjoint operator on $\ell^2:=\ell^2(\{1,2,3,\dots\})$, together with the
Cauchy interlacing theorem and the Courant--Fischer min--max principle.

\section{The infinite Hilbert operator}

Let $\ell^2$ be the Hilbert space of square-summable sequences
$x=(x_1,x_2,\dots)$ with inner product $\langle x,y\rangle=\sum_{k\ge1}x_k\overline{y_k}$ and
norm $\|x\|=\langle x,x\rangle^{1/2}$. Let $\{e_k\}_{k\ge1}$ be the standard orthonormal basis.
Define the infinite Hilbert matrix as the linear operator $H$ determined on this basis by
\begin{equation}\label{eq:Hdef}
\langle He_j,e_i\rangle=\frac{1}{i+j-1},\qquad i,j\ge1 .
\end{equation}

\begin{lemma}[Boundedness, norm, positivity]\label{lem:bounded}
The expression \eqref{eq:Hdef} defines a bounded self-adjoint operator $H$ on $\ell^2$ with
operator norm $\|H\|=\pi$, and $0\le H\le \pi I$ in the sense of quadratic forms. Concretely, for
every finitely supported $x\in\ell^2$,
\begin{equation}\label{eq:hilbertineq}
0\ \le\ \langle Hx,x\rangle=\sum_{i,j\ge1}\frac{x_i\,\overline{x_j}}{i+j-1}\ \le\ \pi\,\|x\|^2 ,
\end{equation}
the constant $\pi$ being best possible.
\end{lemma}

\begin{proof}
The matrix $\bigl(\tfrac1{i+j-1}\bigr)$ is real and symmetric, so the densely defined operator
on finitely supported sequences is symmetric. The upper bound in \eqref{eq:hilbertineq} with the
sharp constant $\pi$ is Hilbert's inequality in its symmetric (``diagonal'') form
$\sum_{i,j}\frac{a_i a_j}{i+j-1}\le\pi\sum_i a_i^2$; see Hardy, Littlewood and P\'olya,
\emph{Inequalities}, 2nd ed., Cambridge Univ.\ Press, 1952, Theorem~315, together with the
discussion of sharpness in \S9.1. This shows $H$ is bounded of norm $\le\pi$ on the dense set of
finitely supported sequences, hence extends uniquely to a bounded operator on $\ell^2$ with
$\|H\|\le\pi$; being the continuous extension of a symmetric operator, it is self-adjoint.

For positivity, use $\frac1{i+j-1}=\int_0^1 t^{\,i+j-2}\,dt$. For finitely supported $x$ the
double sum is finite, so we may interchange sum and integral:
\[
\langle Hx,x\rangle=\sum_{i,j\ge1}x_i\overline{x_j}\int_0^1 t^{\,i+j-2}\,dt
=\int_0^1\Bigl(\sum_{i}x_i t^{\,i-1}\Bigr)\overline{\Bigl(\sum_{j}x_j t^{\,j-1}\Bigr)}\,dt
=\int_0^1\Bigl|\sum_{k\ge1}x_k\,t^{\,k-1}\Bigr|^2\,dt\ \ge\ 0 .
\]
Hence $0\le\langle Hx,x\rangle$ on a dense set, so $H\ge0$, and with the upper bound,
$0\le H\le\pi I$. The constant $\pi$ is best possible by Hardy--Littlewood--P\'olya
(loc.\ cit.), so $\|H\|=\pi$.
\end{proof}

\begin{lemma}[Spectrum]\label{lem:spec}
The spectrum of $H$ is $\spec(H)=[0,\pi]$, and $H$ has \emph{no} eigenvalues; i.e.\ its point
spectrum is empty and the spectrum is purely continuous. Consequently, for the spectral measure
$E(\cdot)$ of $H$ and any Borel set $B\subseteq\mathbb{R}$, the spectral projection $E(B)$ is
either $0$ or of infinite rank.
\end{lemma}

\begin{proof}
The identification $\spec(H)=[0,\pi]$ together with the absence of eigenvalues is the classical
spectral analysis of the Hilbert matrix due to W.~Magnus, \emph{On the spectrum of Hilbert's
matrix}, Amer.\ J.\ Math.\ \textbf{72} (1950), 699--704; see also M.~Rosenblum, \emph{On the
Hilbert matrix I, II}, Proc.\ Amer.\ Math.\ Soc.\ \textbf{9} (1958), 137--140 and 581--585, and
S.~C.~Power, \emph{Hankel Operators on Hilbert Space}, Pitman, 1982, Ch.~3. In that analysis $H$
is realized as a Hankel operator that is unitarily equivalent, via the Mellin transform, to
multiplication by the function $t\mapsto \pi/\cosh(\pi t)$ on $L^2(\mathbb{R})$. A multiplication
operator $M_\varphi$ by a continuous, non-constant, real function $\varphi$ has spectrum equal to
the closure of the range of $\varphi$ and has no eigenvalues (since $\{t:\varphi(t)=\lambda\}$ has
Lebesgue measure zero for every $\lambda$, because $\varphi$ is non-constant and real-analytic);
the range of $\pi/\cosh(\pi t)$ over $t\in\mathbb{R}$ is $(0,\pi]$, whose closure is $[0,\pi]$.
Thus $\spec(H)=[0,\pi]$ and $H$ has empty point spectrum.

It remains to deduce the last sentence. Suppose $E(B)$ is a nonzero spectral projection of
\emph{finite} rank $m\ge1$. Its range $M=\ran E(B)$ is an $m$-dimensional subspace that reduces
$H$ (that is, $H(M)\subseteq M$, because spectral projections commute with $H$), and the
restriction $H|_M$ is a self-adjoint operator on the finite-dimensional space $M$. By the
spectral theorem in finite dimensions, $H|_M$ has an eigenvector $v\in M$, $v\ne0$, with some
eigenvalue $\lambda$; since $M$ reduces $H$, $Hv=\lambda v$ holds as an identity in $\ell^2$, so
$\lambda$ is an eigenvalue of $H$, contradicting the absence of eigenvalues. Hence every nonzero
$E(B)$ has infinite rank.
\end{proof}

\section{Compression to finite sections}

Let $P_n\colon\ell^2\to\ell^2$ be the orthogonal projection onto
$V_n:=\operatorname{span}\{e_1,\dots,e_n\}$, i.e.\ $P_nx=(x_1,\dots,x_n,0,0,\dots)$. Identifying
$V_n$ with $\mathbb{C}^n$ via $e_1,\dots,e_n$, the compression $A_n:=P_nHP_n|_{V_n}$ has matrix
$\bigl(\langle He_j,e_i\rangle\bigr)_{i,j=1}^n=\bigl(\tfrac1{i+j-1}\bigr)_{i,j=1}^n=H_n$. Thus
$\sigma_i(H_n)$ is the $i$-th largest eigenvalue $\lambda_i(A_n)$ of the self-adjoint operator
$A_n$ on $V_n$. Note that for $x\in V_n$ we have $P_nx=x$, hence
\begin{equation}\label{eq:rayleigh}
\langle A_n x,x\rangle=\langle P_nHP_n x,x\rangle=\langle Hx,x\rangle ,\qquad x\in V_n .
\end{equation}

We use the Courant--Fischer min--max principle: for a self-adjoint operator $A$ on a
\emph{finite}-dimensional inner product space $\mathcal H$ and $1\le i\le\dim\mathcal H$,
\begin{equation}\label{eq:minmax}
\lambda_i(A)=\max_{\substack{S\subseteq \mathcal H \\ \dim S=i}}\ \min_{\substack{x\in S\\ x\neq0}}
\frac{\langle Ax,x\rangle}{\|x\|^2}
\end{equation}
(Horn--Johnson, \emph{Matrix Analysis}, 2nd ed., Theorem~4.2.6).

\subsection{Upper bound, strictness, and monotonicity}

\begin{lemma}\label{lem:upper}
For every $n\ge i\ge1$ we have $\sigma_i(H_n)<\pi$. Furthermore
$\sigma_i(H_{n})\le\sigma_i(H_{n+1})$ for all $n\ge i$.
\end{lemma}

\begin{proof}
\textbf{Upper bound $\le\pi$.} Let $x_0\in V_n$ be a unit eigenvector of $A_n$ for the eigenvalue
$\sigma_i(H_n)=\lambda_i(A_n)$. By \eqref{eq:rayleigh} and Lemma~\ref{lem:bounded},
$\sigma_i(H_n)=\langle A_nx_0,x_0\rangle=\langle Hx_0,x_0\rangle\le\pi\|x_0\|^2=\pi$.

\textbf{Strictness.} If $\sigma_i(H_n)=\pi$, then $\sigma_1(H_n)\ge\sigma_i(H_n)=\pi$, so there is
a unit vector $x\in V_n$ with $\langle Hx,x\rangle=\langle A_nx,x\rangle=\pi=\|H\|\|x\|^2$. Since
$\pi I-H\ge0$ and $\langle(\pi I-H)x,x\rangle=0$, positivity of $\pi I-H$ gives
$(\pi I-H)^{1/2}x=0$, hence $(\pi I-H)x=0$, i.e.\ $Hx=\pi x$. Then $\pi$ would be an eigenvalue of
$H$, contradicting Lemma~\ref{lem:spec}. Therefore $\sigma_i(H_n)<\pi$ for every finite $n$.

\textbf{Monotonicity.} $H_n$ is the leading $n\times n$ principal submatrix of $H_{n+1}$,
obtained by deleting the last row and column. By the Cauchy interlacing theorem
(Horn--Johnson, \emph{Matrix Analysis}, 2nd ed., Theorem~4.3.17),
\[
\sigma_{i}(H_{n+1})\ \ge\ \sigma_{i}(H_{n})\ \ge\ \sigma_{i+1}(H_{n+1}),\qquad 1\le i\le n .
\]
The left inequality is the asserted monotonicity.
\end{proof}

\subsection{Lower bound: each fixed singular value approaches $\pi$}

\begin{lemma}\label{lem:lower}
Fix $i\ge1$ and $0<\varepsilon\le\pi$. Then there exists $N=N(i,\varepsilon)$ such that
$\sigma_i(H_n)>\pi-\varepsilon$ for all $n\ge N$.
\end{lemma}

\begin{proof}
Let $E(\cdot)$ be the spectral measure of $H$ and set $B:=(\pi-\tfrac{\varepsilon}{2},\,\pi]$.
Since $\spec(H)=[0,\pi]$ (Lemma~\ref{lem:spec}) and $\pi\in B$, the interval $B$ meets the
spectrum, so $E(B)\ne0$; by Lemma~\ref{lem:spec} its range $W:=\ran E(B)$ is infinite-dimensional.
Choose an orthonormal system $w_1,\dots,w_i\in W$ and put
$U:=\operatorname{span}\{w_1,\dots,w_i\}$, an $i$-dimensional subspace of $\ell^2$ contained in
$W=\ran E(B)$.

Fix once and for all the threshold
\begin{equation}\label{eq:deltachoice}
\delta:=\frac{\varepsilon}{40\,\pi\,(i+1)}\ \in(0,\tfrac1{40}] ,
\end{equation}
where the bound $\delta\le\frac1{40}$ uses $\varepsilon\le\pi$ and $i+1\ge2$.

\smallskip
\noindent\emph{Step 1: the quadratic form is large on $U$.}
For any unit vector $u\in U\subseteq\ran E(B)$ we have $u=E(B)u$, so by the spectral theorem the
scalar measure $\mu_u(\cdot):=\langle E(\cdot)u,u\rangle$ is supported in $B$ and
$\mu_u(B)=\|u\|^2=1$. Since $\lambda\ge\pi-\tfrac\varepsilon2$ for all $\lambda\in B$,
\begin{equation}\label{eq:formlarge}
\langle Hu,u\rangle=\int_{[0,\pi]}\lambda\,d\mu_u(\lambda)=\int_{B}\lambda\,d\mu_u(\lambda)
\ \ge\ \bigl(\pi-\tfrac\varepsilon2\bigr)\,\mu_u(B)=\pi-\tfrac\varepsilon2 .
\end{equation}

\smallskip
\noindent\emph{Step 2: truncate $U$ into a finite section $V_N$.}
Each $w_m\in\ell^2$, so $\sum_{k>N}|\langle w_m,e_k\rangle|^2\to0$ as $N\to\infty$. Choose $N$ so
large that
\begin{equation}\label{eq:tail}
\sum_{k>N}|\langle w_m,e_k\rangle|^2<\delta^2\qquad(m=1,\dots,i) ,
\end{equation}
and set $w_m':=P_N w_m\in V_N$. By \eqref{eq:tail},
$\|w_m-w_m'\|=\bigl(\sum_{k>N}|\langle w_m,e_k\rangle|^2\bigr)^{1/2}<\delta$ for each $m$, and
$\|w_m'\|\le\|w_m\|=1$. Let $U':=\operatorname{span}\{w_1',\dots,w_i'\}\subseteq V_N$.

\smallskip
\noindent\emph{Step 2a: $\dim U'=i$.} Let
$G:=\bigl(\langle w_m',w_\ell'\rangle\bigr)_{m,\ell=1}^i$ be the Gram matrix of the $w_m'$, and
recall the Gram matrix of the orthonormal $w_m$ is $I_i$. For each $m,\ell$,
\[
\bigl|\langle w_m',w_\ell'\rangle-\langle w_m,w_\ell\rangle\bigr|
\le\bigl|\langle w_m'-w_m,\,w_\ell'\rangle\bigr|+\bigl|\langle w_m,\,w_\ell'-w_\ell\rangle\bigr|
\le\delta\,\|w_\ell'\|+\delta\,\|w_m\|\le2\delta ,
\]
so $\|G-I_i\|_{\mathrm{op}}\le\|G-I_i\|_{\mathrm F}\le2i\delta$. By \eqref{eq:deltachoice},
$2i\delta\le2(i+1)\delta=\frac{\varepsilon}{20\pi}\le\frac1{20}<1$, hence $G$ is positive
definite and invertible, so $w_1',\dots,w_i'$ are linearly independent and $\dim U'=i$.

\smallskip
\noindent\emph{Step 2b: Rayleigh quotient on $U'$.} Let $y=\sum_{m=1}^i c_m w_m'\in U'$ be an
arbitrary unit vector, $\|y\|=1$, with coefficient vector $c=(c_1,\dots,c_i)$. From
$\|G-I_i\|_{\mathrm{op}}\le2i\delta\le\frac1{20}$ we get
$1=\|y\|^2=c^*Gc\ge(1-2i\delta)\|c\|_2^2\ge\tfrac{19}{20}\|c\|_2^2$, so
$\|c\|_2\le\sqrt{20/19}\le2$. Set $z:=\sum_{m=1}^i c_m w_m\in U$. Then
\[
\|y-z\|=\Bigl\|\sum_{m=1}^i c_m(w_m'-w_m)\Bigr\|\le\|c\|_2\,\Bigl(\sum_{m}\|w_m'-w_m\|^2\Bigr)^{1/2}
\le2\cdot\sqrt{i}\,\delta=2\sqrt i\,\delta ,
\]
where the middle step is the Cauchy--Schwarz/triangle estimate
$\|\sum c_m v_m\|\le\|c\|_2(\sum\|v_m\|^2)^{1/2}$ applied to $v_m=w_m'-w_m$. By
\eqref{eq:deltachoice}, $2\sqrt i\,\delta\le2(i+1)\delta=\frac{\varepsilon}{20\pi}\le\frac1{20}$,
hence
\begin{equation}\label{eq:yz}
\|y-z\|\le2\sqrt i\,\delta\le\tfrac1{20},\qquad
\|z\|\ge\|y\|-\|y-z\|\ge1-\tfrac1{20}=\tfrac{19}{20}>\tfrac12 ,\qquad
\|z\|\le\|y\|+\|y-z\|\le2 .
\end{equation}

Now, since $y\in V_N$, \eqref{eq:rayleigh} gives $\langle A_Ny,y\rangle=\langle Hy,y\rangle$.
Estimate the difference $\langle Hy,y\rangle-\langle Hz,z\rangle$ using bilinearity and
$\|H\|=\pi$:
\[
\bigl|\langle Hy,y\rangle-\langle Hz,z\rangle\bigr|
=\bigl|\langle H(y-z),y\rangle+\langle Hz,\,y-z\rangle\bigr|
\le\pi\|y-z\|\,\|y\|+\pi\|z\|\,\|y-z\|
\le\pi(1+2)\,\|y-z\| ,
\]
using $\|y\|=1$ and $\|z\|\le2$ from \eqref{eq:yz}; thus
$\bigl|\langle Hy,y\rangle-\langle Hz,z\rangle\bigr|\le3\pi\|y-z\|\le6\pi\sqrt i\,\delta$.
Since $z\in U$ with $\|z\|>\tfrac12$, the unit vector $z/\|z\|$ lies in $U$, so by
\eqref{eq:formlarge},
\[
\langle Hz,z\rangle=\|z\|^2\Bigl\langle H\tfrac{z}{\|z\|},\tfrac{z}{\|z\|}\Bigr\rangle
\ \ge\ \|z\|^2\bigl(\pi-\tfrac\varepsilon2\bigr) .
\]
By \eqref{eq:yz}, $\|z\|^2\ge(1-2\sqrt i\,\delta)^2\ge1-4\sqrt i\,\delta$, and since
$0\le\pi-\tfrac\varepsilon2\le\pi$,
\[
\|z\|^2\bigl(\pi-\tfrac\varepsilon2\bigr)\ \ge\ \bigl(\pi-\tfrac\varepsilon2\bigr)-4\sqrt i\,\delta\cdot\pi .
\]
Combining the last three displays,
\[
\langle A_Ny,y\rangle=\langle Hy,y\rangle
\ \ge\ \langle Hz,z\rangle-6\pi\sqrt i\,\delta
\ \ge\ \bigl(\pi-\tfrac\varepsilon2\bigr)-4\pi\sqrt i\,\delta-6\pi\sqrt i\,\delta
=\pi-\tfrac\varepsilon2-10\pi\sqrt i\,\delta .
\]
Finally, by \eqref{eq:deltachoice} and $\sqrt i\le i+1$,
$10\pi\sqrt i\,\delta\le10\pi(i+1)\delta=\tfrac{\varepsilon}{4}$, so for every unit vector
$y\in U'$,
\begin{equation}\label{eq:Uprimeform}
\langle A_Ny,y\rangle\ \ge\ \pi-\tfrac\varepsilon2-\tfrac\varepsilon4\ \ge\ \pi-\varepsilon .
\end{equation}

\smallskip
\noindent\emph{Step 3: conclude via min--max.} The subspace $U'\subseteq V_N$ has dimension $i$
(Step~2a), and by \eqref{eq:Uprimeform} the Rayleigh quotient of $A_N$ exceeds $\pi-\varepsilon$
on $U'\setminus\{0\}$. By the min--max principle \eqref{eq:minmax} for $A_N$ on $V_N$,
\[
\sigma_i(H_N)=\lambda_i(A_N)
=\max_{\substack{S\subseteq V_N\\\dim S=i}}\ \min_{\substack{y\in S\\ y\ne0}}
\frac{\langle A_Ny,y\rangle}{\|y\|^2}
\ \ge\ \min_{\substack{y\in U'\\ y\ne0}}\frac{\langle A_Ny,y\rangle}{\|y\|^2}\ \ge\ \pi-\varepsilon .
\]
By the monotonicity of Lemma~\ref{lem:upper}, $\sigma_i(H_n)\ge\sigma_i(H_N)\ge\pi-\varepsilon$
for all $n\ge N$.
\end{proof}

\section{Proof of Theorem~\ref{thm:main}}

\begin{proof}
Fix $i\ge1$. By Lemma~\ref{lem:upper}, the sequence $\bigl(\sigma_i(H_n)\bigr)_{n\ge i}$ is
nondecreasing and bounded above by $\pi$ (strictly, $\sigma_i(H_n)<\pi$), so the limit
$L_i:=\lim_{n\to\infty}\sigma_i(H_n)$ exists and $L_i\le\pi$. By Lemma~\ref{lem:lower}, for every
$\varepsilon\in(0,\pi]$ there is $N$ with $\sigma_i(H_n)>\pi-\varepsilon$ for all $n\ge N$, hence
$L_i\ge\pi-\varepsilon$. Letting $\varepsilon\downarrow0$ gives $L_i\ge\pi$. Therefore $L_i=\pi$:
\[
\lim_{n\to\infty}\sigma_i(H_n)=\pi\qquad\text{for every fixed }i\ge1 .
\]
The strict inequality $\sigma_i(H_n)<\pi$ for all finite $n$ (Lemma~\ref{lem:upper}) shows the
limit is approached but never attained.
\end{proof}

\section{Rate of convergence}

The equality $\lim_n\sigma_i(H_n)=\pi$ holds, so the proposed statement is \emph{true} for every
fixed $i$; the convergence is only of order $1/\log n$. We record the precise rate for
completeness; it is not used in the proof of Theorem~\ref{thm:main}.

\begin{proposition}\label{prop:rate}
For each fixed $i\ge1$, $\;0<\pi-\sigma_i(H_n)=O\!\left(\dfrac1{\log n}\right)$ as $n\to\infty$.
For the largest eigenvalue the sharp leading asymptotics are the de Bruijn--Wilf relation
\[
\pi-\sigma_1(H_n)\ \sim\ \frac{\pi^3}{\log n}\qquad(n\to\infty),
\]
and $\pi-\sigma_i(H_n)=\Theta(1/\log n)$ for every fixed $i$.
\end{proposition}

\begin{remark}
The asymptotic $\pi-\sigma_1(H_n)\sim\pi^3/\log n$ is the theorem of N.~G.~de Bruijn and
H.~S.~Wilf, \emph{On Hilbert's inequality in $n$ dimensions}, Bull.\ Amer.\ Math.\ Soc.\
\textbf{68} (1962), 70--73; finer corrections (lower-order terms involving $\log\log n$) appear in
the subsequent literature. Direct numerical evaluation of $\sigma_1(H_n)$ for $n$ up to several
thousand gives values of $(\pi-\sigma_1(H_n))\log n$ around $4$--$5$, still well below
$\pi^3\approx31.0$; this reflects the large $\log\log n/\log n$ correction at such $n$ and is
consistent with the de Bruijn--Wilf expansion. Theorem~\ref{thm:main} does not rely on
Proposition~\ref{prop:rate}: its only quantitative inputs are the qualitative spectral facts of
Lemmas~\ref{lem:bounded}--\ref{lem:spec}.
\end{remark}

\end{document}
